% Part: first-order-logic
% Chapter: completeness
% Section: complete-sets

% Definition of complete consistent sets. Properties of complete sets required
% for completeness proved are in provability.tex in the
% chapter on the proof system used.

\documentclass[../../../include/open-logic-section]{subfiles}

\begin{document}

\iftag{FOL}
      {\olfileid{fol}{com}{ccs}}
      {\olfileid{pl}{com}{ccs}}
      
\olsection{Volledige en consistente verzamelingen van \usetoken{P}{sentence}}

\begin{defn}[Volledige verzameling]
\ollabel{def:complete-set} Een verzameling~$\Gamma$ van !!{sentence}s heet
\emph{!!{complete}} precies als voor elke !!{sentence}~$!A$ minstens een
van deze twee mogelijkheden geldt: $!A \in \Gamma$ of
$\lnot !A \in \Gamma$.
\end{defn}

\begin{explain}
Een !!^{complete} verzameling zinnen laat geen enkele vraag open. Neem
een willekeurige !!{sentence}~$!A$. De verzameling~$\Gamma$ ``zegt'' dan
of $!A$ waar of onwaar is. Zulke verzamelingen zijn niet alleen nodig
voor het bewijs van de volledigheidsstelling. Een theorie die !!{complete}
is en met axioma's kan worden vastgelegd, is bijvoorbeeld altijd
beslisbaar.
\end{explain}

\begin{explain}
In het volledigheidsbewijs gebruiken we verzamelingen die zowel
!!{complete} als consistent zijn. Elke consistente verzameling
!!{sentence}s~$\Gamma$ kan namelijk worden opgenomen in een grotere
verzameling~$\Gamma^*$ die volledig én consistent is. Voor elke
!!{sentence}~$!A$ bevat zo'n verzameling ofwel $!A$ ofwel haar ontkenning
$\lnot !A$, maar niet beide. Dit geldt dus ook voor
\iftag{FOL}{atomaire !!{sentence}s}{propositionele variabelen}. Uit
!!a{complete} consistente verzameling\iftag{FOL}{ in een taal die we op
de juiste manier hebben uitgebreid met !!{constant}s}{} kunnen we daarom
\iftag{FOL}{!!a{structure}}{!!a{valuation}} bouwen. Welke
\iftag{FOL}{betekenis de !!{predicate}s krijgen}{waarheidswaarde de
propositionele variabelen krijgen}, laten we afhangen van de
\iftag{FOL}{atomaire !!{sentence}s}{!!{propositional variable}s} in
~$\Gamma^*$. Daarna bewijzen we dat deze
\iftag{FOL}{!!{structure}}{!!{valuation}} alle !!{sentence}s in
~$\Gamma^*$ waar maakt, en daarmee ook alle zinnen in~$\Gamma$. Voor dat
bewijs hebben we onder meer nodig dat $\lnot !A \in \Gamma^*$ precies
als $!A \notin \Gamma^*$, en dat $(!A \lor !B) \in \Gamma^*$ precies
als $!A \in \Gamma^*$ of $!B \in \Gamma^*$. Voor de andere
verbindingswoorden hebben we soortgelijke regels nodig.
\end{explain}

Hierna gebruiken we vaak zonder het er telkens bij te zeggen dat
$\Proves$ reflexief, monotoon en transitief is (zie
\iftag{FOL}{%
  \tagrefs{prfSC/{fol:seq:ptn:sec},prfND/{fol:ntd:ptn:sec},prfAX/{fol:axd:ptn:sec},prfTab/{fol:tab:ptn:sec}}}{%
  \tagrefs{prfSC/{pl:seq:ptn:sec},prfND/{pl:ntd:ptn:sec},prfAX/{pl:axd:ptn:sec},prfTab/{pl:tab:ptn:sec}}}).

\begin{prop}
\ollabel{prop:ccs}
Stel dat $\Gamma$ !!{complete} en consistent is. Dan geldt:
\begin{enumerate}
\item \ollabel{prop:ccs-prov-in} Als $\Gamma \Proves !A$, dan
  $!A \in \Gamma$.

\tagitem{prvAnd}{\ollabel{prop:ccs-and} $!A \land !B \in \Gamma$
  precies als zowel $!A \in \Gamma$ als $!B \in \Gamma$.}{}

\tagitem{prvOr}{\ollabel{prop:ccs-or} $!A \lor !B \in \Gamma$ precies
  als $!A \in \Gamma$ of $!B \in \Gamma$.}{}

\tagitem{prvIf}{\ollabel{prop:ccs-if} $!A \lif !B \in \Gamma$ precies
  als $!A \notin \Gamma$ of $!B \in \Gamma$.}{}
\end{enumerate}
\end{prop}

\begin{proof}
In de rest van dit bewijs nemen we steeds aan dat $\Gamma$ !!{complete}
en consistent is.
\begin{enumerate}
\item Als $\Gamma \Proves !A$, dan $!A \in \Gamma$.

Stel dat $\Gamma \Proves !A$. Neem nu aan dat $!A \notin \Gamma$.
Omdat $\Gamma$ !!{complete} is, geldt dan $\lnot !A \in \Gamma$. Volgens
\iftag{FOL}{%
  \tagrefs{prfSC/{fol:seq:prv:prop:explicit-inc},
    prfND/{fol:ntd:prv:prop:explicit-inc},
    prfAX/{fol:axd:prv:prop:explicit-inc},
    prfTab/{fol:tab:prv:prop:explicit-inc}}}{%
  \tagrefs{prfSC/{pl:seq:prv:prop:explicit-inc},
    prfND/{pl:ntd:prv:prop:explicit-inc},
    prfAX/{pl:axd:prv:prop:explicit-inc},
    prfTab/{pl:tab:prv:prop:explicit-inc}}}
zou $\Gamma$ inconsistent zijn. Maar we hebben juist aangenomen dat
$\Gamma$ consistent is. Onze extra aanname kan dus niet kloppen:
$!A \notin \Gamma$ is onmogelijk, zodat $!A \in \Gamma$.

\tagitem{defAnd}{}{%
\iftag{probAnd}{Opgave.}{%
$!A \land !B \in \Gamma$ precies als zowel $!A \in \Gamma$ als
$!B \in \Gamma$:

Eerst van links naar rechts. Stel dat $!A \land !B \in \Gamma$. Uit
\iftag{FOL}{%
  \tagrefs{prfSC/{fol:seq:ppr:prop:provability-land},
    prfND/{fol:ntd:ppr:prop:provability-land},
    prfAX/{fol:axd:ppr:prop:provability-land},
    prfTab/{fol:tab:ppr:prop:provability-land}}}{%
  \tagrefs{prfSC/{pl:seq:ppr:prop:provability-land},
    prfND/{pl:ntd:ppr:prop:provability-land},
    prfAX/{pl:axd:ppr:prop:provability-land},
    prfTab/{pl:tab:ppr:prop:provability-land}}}, onderdeel~(1),
volgt dan $\Gamma \Proves !A$ en $\Gamma \Proves !B$. Volgens
\olref{prop:ccs-prov-in} geldt dus $!A \in \Gamma$ en
$!B \in \Gamma$. Dat is wat we moesten aantonen.

Nu van rechts naar links. Neem $!A \in \Gamma$ en $!B \in \Gamma$
aan. Uit
\iftag{FOL}{%
  \tagrefs{prfSC/{fol:seq:ppr:prop:provability-land},
    prfND/{fol:ntd:ppr:prop:provability-land},
    prfAX/{fol:axd:ppr:prop:provability-land},
    prfTab/{fol:tab:ppr:prop:provability-land}}}{%
  \tagrefs{prfSC/{pl:seq:ppr:prop:provability-land},
    prfND/{pl:ntd:ppr:prop:provability-land},
    prfAX/{pl:axd:ppr:prop:provability-land},
    prfTab/{pl:tab:ppr:prop:provability-land}}}, onderdeel~(2),
volgt $\Gamma \Proves !A \land !B$. Volgens
\olref{prop:ccs-prov-in} geldt dus $!A \land !B \in \Gamma$.}}

\tagitem{defOr}{}{%
\iftag{probOr}{Opgave.}{%
Eerst bewijzen we: als $!A \lor !B \in \Gamma$, dan geldt
$!A \in \Gamma$ of $!B \in \Gamma$. Stel dat
$!A \lor !B \in \Gamma$, maar dat $!A \notin \Gamma$ en
$!B \notin \Gamma$. Omdat $\Gamma$ !!{complete} is, geldt
$\lnot !A \in \Gamma$ en $\lnot !B \in \Gamma$. Uit
\iftag{FOL}{%
  \tagrefs{prfSC/{fol:seq:ppr:prop:provability-lor},
    prfND/{fol:ntd:ppr:prop:provability-lor},
    prfAX/{fol:axd:ppr:prop:provability-lor},
    prfTab/{fol:tab:ppr:prop:provability-lor}}}{%
  \tagrefs{prfSC/{pl:seq:ppr:prop:provability-lor},
    prfND/{pl:ntd:ppr:prop:provability-lor},
    prfAX/{pl:axd:ppr:prop:provability-lor},
    prfTab/{pl:tab:ppr:prop:provability-lor}}},
onderdeel~(1), volgt dan dat $\Gamma$ inconsistent is. Dat kan niet.
Dus geldt $!A \in \Gamma$ of $!B \in \Gamma$.

Voor de andere richting nemen we aan dat $!A \in \Gamma$ of
$!B \in \Gamma$. Uit
\iftag{FOL}{%
  \tagrefs{prfSC/{fol:seq:ppr:prop:provability-lor},
    prfND/{fol:ntd:ppr:prop:provability-lor},
    prfAX/{fol:axd:ppr:prop:provability-lor},
    prfTab/{fol:tab:ppr:prop:provability-lor}}}{%
  \tagrefs{prfSC/{pl:seq:ppr:prop:provability-lor},
    prfND/{pl:ntd:ppr:prop:provability-lor},
    prfAX/{pl:axd:ppr:prop:provability-lor},
    prfTab/{pl:tab:ppr:prop:provability-lor}}}, onderdeel~(2),
volgt $\Gamma \Proves !A \lor !B$. Volgens
\olref{prop:ccs-prov-in} geldt dus $!A \lor !B \in \Gamma$.}}

\tagitem{defIf}{}{%
\iftag{probIf}{Opgave.}{%
Eerst van links naar rechts. Stel dat $!A \lif !B \in \Gamma$. Neem
nu voor een tegenspraak aan dat $!A \in \Gamma$ en
$!B \notin \Gamma$. We hebben dan $!A \lif !B \in \Gamma$ en
$!A \in \Gamma$. Uit
\iftag{FOL}{%
  \tagrefs{prfSC/{fol:seq:ppr:prop:provability-lif},
    prfND/{fol:ntd:ppr:prop:provability-lif},
    prfAX/{fol:axd:ppr:prop:provability-lif},
    prfTab/{fol:tab:ppr:prop:provability-lif}}}{%
  \tagrefs{prfSC/{pl:seq:ppr:prop:provability-lif},
    prfND/{pl:ntd:ppr:prop:provability-lif},
    prfAX/{pl:axd:ppr:prop:provability-lif},
    prfTab/{pl:tab:ppr:prop:provability-lif}}}, onderdeel~(1),
volgt $\Gamma \Proves !B$. Volgens \olref{prop:ccs-prov-in} geldt dan
$!B \in \Gamma$. Dat botst met onze aanname dat $!B \notin \Gamma$.

Nu van rechts naar links. Bekijk eerst het geval $!A \notin \Gamma$.
Omdat $\Gamma$ !!{complete} is, geldt dan $\lnot !A \in \Gamma$. Uit
\iftag{FOL}{%
  \tagrefs{prfSC/{fol:seq:ppr:prop:provability-lif},
    prfND/{fol:ntd:ppr:prop:provability-lif},
    prfAX/{fol:axd:ppr:prop:provability-lif},
    prfTab/{fol:tab:ppr:prop:provability-lif}}}{%
  \tagrefs{prfSC/{pl:seq:ppr:prop:provability-lif},
    prfND/{pl:ntd:ppr:prop:provability-lif},
    prfAX/{pl:axd:ppr:prop:provability-lif},
    prfTab/{pl:tab:ppr:prop:provability-lif}}}, onderdeel~(2),
volgt $\Gamma \Proves !A \lif !B$. Volgens
\olref{prop:ccs-prov-in} geldt dus $!A \lif !B \in \Gamma$.

Bekijk ten slotte het geval $!B \in \Gamma$. Uit
\iftag{FOL}{%
  \tagrefs{prfSC/{fol:seq:ppr:prop:provability-lif},
    prfND/{fol:ntd:ppr:prop:provability-lif},
    prfTab/{fol:tab:ppr:prop:provability-lif},
    prfAX/{fol:axd:ppr:prop:provability-lif}}}{%
  \tagrefs{prfSC/{pl:seq:ppr:prop:provability-lif},
    prfND/{pl:ntd:ppr:prop:provability-lif},
    prfTab/{pl:tab:ppr:prop:provability-lif},
    prfAX/{pl:axd:ppr:prop:provability-lif}}},
opnieuw onderdeel~(2), volgt $\Gamma \Proves !A \lif !B$. Volgens
\olref{prop:ccs-prov-in} geldt dus $!A \lif !B \in \Gamma$.}}
\end{enumerate}
\end{proof}

\tagprob{FOL}
\begin{prob}
Maak het bewijs van \olref[fol][com][ccs]{prop:ccs} af.
\end{prob}
\tagendprob

\tagprob{notFOL}
\begin{prob}
Maak het bewijs van \olref[pl][com][ccs]{prop:ccs} af.
\end{prob}
\tagendprob

\end{document}
